\documentclass{../base/canon}

%% canon_variables.tex — Valores por defecto de metadatos institucionales
%%
%% Incluir ANTES del \begin{document} para sobreescribir los defaults de la clase.
%% El template puede sobreescribir cualquier valor después de este \input.
%%
%% Uso en template:
%%   %% canon_variables.tex — Valores por defecto de metadatos institucionales
%%
%% Incluir ANTES del \begin{document} para sobreescribir los defaults de la clase.
%% El template puede sobreescribir cualquier valor después de este \input.
%%
%% Uso en template:
%%   %% canon_variables.tex — Valores por defecto de metadatos institucionales
%%
%% Incluir ANTES del \begin{document} para sobreescribir los defaults de la clase.
%% El template puede sobreescribir cualquier valor después de este \input.
%%
%% Uso en template:
%%   \input{../base/canon_variables}
%%   \SetDocTitle{Mi Informe Específico}   % sobreescribe el default

\SetDocLogo{../assets/logo.pdf}
\SetDocUnit{Tecnología \& Sistemas}
\SetContactUnit{Mesa de soporte}
\SetContactMembers{%
  Equipo CoreX & Soporte & soporte@empresa.com%
}
\SetContactNote{Para soporte técnico o consultas sobre este documento.}

%%   \SetDocTitle{Mi Informe Específico}   % sobreescribe el default

\SetDocLogo{../assets/logo.pdf}
\SetDocUnit{Tecnología \& Sistemas}
\SetContactUnit{Mesa de soporte}
\SetContactMembers{%
  Equipo CoreX & Soporte & soporte@empresa.com%
}
\SetContactNote{Para soporte técnico o consultas sobre este documento.}

%%   \SetDocTitle{Mi Informe Específico}   % sobreescribe el default

\SetDocLogo{../assets/logo.pdf}
\SetDocUnit{Tecnología \& Sistemas}
\SetContactUnit{Mesa de soporte}
\SetContactMembers{%
  Equipo CoreX & Soporte & soporte@empresa.com%
}
\SetContactNote{Para soporte técnico o consultas sobre este documento.}

\SetDocType{REPORTE TÉCNICO}
\SetDocTitle{Verificación de Integridad de Datos ERP — Ciclo Semana 17}
\SetDocCode{REP-TEC-003}
\SetDocArea{Tecnología \& Sistemas}
\SetDocAuthor{Equipo de Integración}
\SetDocDate{2026-04-22}
\SetContactUnit{Ingeniería de Datos}
\SetContactMembers{%
  Diego Arévalo  & Ingeniero de Datos & d.arevalo@empresa.com \\
  Ana Viteri     & DBA                & a.viteri@empresa.com%
}

\begin{document}

\section{Contexto}

\begin{ParrafoMetodologico}
Verificación ejecutada el 2026-04-22 sobre la base de datos de producción
(\texttt{corex\_prod}) en el esquema \texttt{postcosecha}.
Herramientas: psql 15.4, script \CodePath{scripts/validate\_week\_integrity.sql}.
Período analizado: semana 17-2026 (2026-04-14 al 2026-04-20).
\end{ParrafoMetodologico}

\section{Hallazgos}

\begin{KeyFindingBox}[Integridad confirmada en BAL1 y BAL2]
  Los totales de BAL1 y BAL2 concilian con el registro de nómina al 99.97\,\%.
  La diferencia de 3 tallos se atribuye a un registro manual corregido post-turno,
  documentado en el log de auditoría \texttt{AUD-BAL-2026-17-003}.
\end{KeyFindingBox}

\begin{WarningBox}[Dato faltante — BAL2 madrugada 17-abr]
  El servidor de adquisición BAL2 no registró datos entre las 02:15 y las 04:40 del
  2026-04-17 por mantenimiento no programado.  El gap afecta 88 registros.
  Se aplicó imputación por promedio de los 3 días anteriores (método \texttt{imp\_avg3}).
\end{WarningBox}

\section{Evidencia técnica}

\subsection{Consulta de reconciliación}

\begin{CodeBlock}
-- Reconciliación BAL1 vs Nómina — Semana 17
SELECT
    b.semana,
    SUM(b.tallos_salida)         AS tallos_bal1,
    SUM(n.horas_reportadas * 85) AS estimado_nomina,
    ABS(SUM(b.tallos_salida)
      - SUM(n.horas_reportadas * 85)) AS diferencia
FROM balanzas_semana b
JOIN nomina_turno n
    ON b.turno_id = n.turno_id
    AND b.semana  = n.semana
WHERE b.semana = 17 AND b.anio = 2026
GROUP BY b.semana;
\end{CodeBlock}

\subsection{Resultado}

\begin{table}[H]
  \centering
  \caption{Resultado de reconciliación BAL1 vs Nómina}
  \begin{tabular}{l r r r r}
    \toprule
    \textbf{Fuente} & \textbf{Tallos} & \textbf{Est. Nómina} & \textbf{Diferencia} & \textbf{Var.\,\%} \\
    \midrule
    BAL1 Semana 17 & 271\,300 & 271\,297 & 3 & 0.001 \\
    \midrule
    \multicolumn{3}{l}{Umbral de tolerancia} & $\leq 500$ & $\leq 0.18\,\%$ \\
    \bottomrule
  \end{tabular}
\end{table}

\subsection{Registros con imputación}

\begin{CodeBlock}
-- Registros imputados en el gap BAL2
SELECT id, timestamp_utc, tallos_raw, tallos_imputado, metodo_imputacion
FROM balanzas_raw
WHERE semana = 17 AND anio = 2026
  AND metodo_imputacion IS NOT NULL
ORDER BY timestamp_utc;
-- 88 filas retornadas
\end{CodeBlock}

\section{Recomendaciones técnicas}

\begin{enumerate}
  \item Instalar UPS en el servidor de adquisición BAL2 para prevenir gaps por
        cortes de energía — \textit{Infraestructura, antes de semana 20}.
  \item Configurar alerta automática cuando no se reciban datos de BAL2 por más
        de 15 minutos en horario operativo — \textit{Monitoreo, plazo: 2026-04-30}.
  \item Documentar el método \texttt{imp\_avg3} en el manual de operaciones del
        sistema de balanzas — \textit{Equipo de Datos, plazo: 2026-05-07}.
\end{enumerate}

\SignatureBlock{Diego Arévalo}{Ingeniero de Datos Senior}

\ContactSignature

\end{document}
